\documentclass{article}
\usepackage[utf8]{inputenc}
\usepackage{amsmath}
\usepackage{geometry}
\geometry{margin=2cm}
\begin{document}

\section*{Análise da Questão 166 - ENEM 2024 - Matemática}

\subsection*{Solução Técnica}

A densidade demográfica $d$ é dada pela razão:

\[
d = \frac{\text{número de habitantes}}{\text{área}},
\]

sendo expressa em habitantes por quilômetro quadrado.

Seja $R$ a região principal e $Q$ uma de suas subdivisões. Dados:
\begin{itemize}
  \item $A_Q = \frac{3}{4} A_R$
  \item $H_Q = \frac{1}{2} H_R$
  \item Densidades: $d(R) = \frac{H_R}{A_R}$ e $d(Q) = \frac{H_Q}{A_Q}$
\end{itemize}

Substituindo, temos:

\[
\begin{aligned}
  d(Q) &= \frac{H_Q}{A_Q} = \frac{\frac{1}{2} H_R}{\frac{3}{4} A_R} = \frac{1}{2} H_R \times \frac{4}{3 A_R} = \frac{2}{3} \frac{H_R}{A_R} = \frac{2}{3} d(R).
\end{aligned}
\]

Portanto, a alternativa correta é:

\[
d(Q) = \frac{2}{3} d(R).
\]

\subsection*{Explicação Didática}

A densidade representa como os habitantes estão distribuídos na área. Embora $Q$ tenha metade dos habitantes de $R$, sua área também é menor, ou seja, $3/4$ da área de $R$. Ao calcular a densidade, devemos considerar ambas as proporções, o que resulta em $\frac{2}{3} d(R)$.

\subsection*{Passo a passo para o aluno}
\begin{enumerate}
  \item Determine as densidades de $R$ e $Q$ com as fórmulas dadas.
  \item Substitua as relações fornecidas para moradores e área.
  \item Simplifique frações cuidadosamente para evitar erros.
  \item Relacione a densidade de $Q$ com a de $R$ para concluir.
\end{enumerate}

\subsection*{Análise das Distratores}

\begin{itemize}
  \item Opção A ($\frac{1}{4} d(R)$): Confusão nas frações sem considerar a relação entre habitantes e área.
  \item Opção B ($\frac{1}{2} d(R)$): Ignorar efeito da área menor na densidade.
  \item Opção C ($\frac{3}{4} d(R)$): Considerar somente a proporção da área.
  \item Opção D ($\frac{3}{2} d(R)$): Erro de inversão na manipulação das frações.
\end{itemize}

\subsection*{Perfil da Questão}

Este problema exige do aluno habilidades de raciocínio com razões e proporções, além de manipulação algébrica básica para identificar corretamente a relação entre densidades dadas as mudanças em população e área.

\subsection*{Conclusão}

A questão apresenta coerência em seu enunciado e solução, exigindo análise crítica da densidade como função da relação entre habitantes e área. A alternativa E é correta e representa a resposta matemática exata. Distratores contemplam erros comuns que facilitam a identificação e compreensão do problema.

\end{document}